\scanpage{30}
et on trouve donc un hom de groupes
\[
\begin{aligned}
&\deleted{\unsure{\tau\longmapsto\varphi_n(\tau)}}\\
&\check{\mathbb V}(L)\xrightarrow{\varphi_n}\operatorname{Aff}\Bigl(\check{\mathbb V}\bigl(\sum_{1\le i\le n}L^{\otimes i}
\Bigr)
\end{aligned}
\]
et un hom
\[
\left\{\begin{aligned}
s&\longmapsto\Sigma^n(\alpha_s)\\
E&\longrightarrow\operatorname{Isom}\Bigl(\check{\mathbb V}\bigl(\sum_{1\le i\le n}L^{\otimes i}\bigr),\unsure{\}|\varphi}\Sigma^n(E)\Bigr)
\end{aligned}\right.
\]
induisant un \deleted{\unsure{isom hom|isomorphisme}} isom
\[
 E\mathop{\wedge}_{\check{\mathbb V}(L)}\check{\mathbb V}\bigl(\sum_{1\le i\le n}L^{\otimes i}\bigr)
 \simeq\Sigma^n(E) \tag*{(34)}
\]
permettant de transporter sur $\Sigma^n(E)$ la
structure affine de $\check{\mathbb V}(\sum_{1\le i\le n}L^{\otimes i})$.

En particulier, on trouve
\[
\varphi_2(\tau)(\sigma_1,\sigma_2)=(\sigma_1+2\tau,\sigma_2+\tau\sigma_1+\tau^2) \tag*{(35)}
\]
la transformation linéaire \deleted{affine} correspondante est
\[
\varphi_2^0(\tau)(\sigma_1,\sigma_2)=(\sigma_1,\sigma_2+\tau\sigma_1) \tag*{(36)}
\]
On \deleted{va} interprête un $s\in\Gamma L$ comme un

\deleted{qui est associé}

\deleted{On peut considérer l'action de $\operatorname{trans}(\tau)$ sur
$F$ donc sur son enveloppe vectorielle \inserted{(\unsure{notée|notation})} $\check B$ \unsure{g ou|ou} $\unsure{\underline{\mathcal O}|\underline{\mathcal E}|\mathfrak E}$),
qui est une extension de \deleted{$\unsure{E|\check E}$} $\underline O_S$ par $L$, on
trouve par tensorisation \inserted{avec $\mathrm{id}_L$} une \deleted{\unsure{exten}} action sur
$L\otimes\check B$, extension de $L$ par $L^{\otimes2}$, \deleted{\unsure{d'où}} soit}

\[
\deleted{\psi(\tau)}\qquad
\deleted{\psi(\tau)\in\deleted{\unsure{\operatorname{AutExt}(}}\operatorname{Autom}\ \text{d'ext}\,(L\otimes\check B)}
\tag*{\deleted{(37)}}
\]

% archive sheet number: 29

